\documentclass[../../../include/open-logic-section]{subfiles}

\begin{document}

\olfileid{sth}{ord-arithmetic}{add}
\olsection{ક્રમવાચક સંખ્યાઓનો સરવાળો}

ધારો કે આપણે $\alpha$ અને $\beta$ ઉમેરવા માગીએ છીએ.
$\alpha$ની એક નકલ પછી તરત $\beta$ની એક નકલ
મૂકી શકીએ. (\emph{નકલો} લેવી પડે, કારણ કે
\olref[ordinals][basic]{ordinalsaresubsets} મુજબ કાં તો
$\alpha
\subseteq \beta$ અથવા $\beta \subseteq \alpha$.)
તેનો સહજ અર્થ એ છે કે પહેલાં $\alpha$-અનુક્રમના
પગલાંમાંથી પસાર થઈએ, અને \emph{ત્યાર પછી}
$\beta$-અનુક્રમમાંથી. મળતો અનુક્રમ સુક્રમિત હશે; તેથી
\olref[ordinals][ordtype]{thmOrdinalRepresentation}
મુજબ તે એક (અનન્ય) ક્રમવાચક સંખ્યા સાથે ક્રમ-એકરૂપ
છે. એ સંખ્યાને $\alpha$ અને $\beta$નો \emph{સરવાળો}
ગણી શકાય.

ક્રમવાચક સરવાળા પાછળનો સહજ વિચાર આ છે. તેની કડક
વ્યાખ્યા આપવા પહેલાં ગણોની \emph{નકલો} લેવાનો વિચાર
જોઈએ. કઈ વસ્તુ ક્યાંથી આવી તેનો હિસાબ રાખવા માટે
$0$ અને $1$ જેવા મનસ્વી ચિહ્નો વાપરીએ છીએ:

\begin{defn}\ollabel{defdissum}
$A$ અને $B$નો \emph{વિચ્છિન્ન સરવાળો} એ
$A \disjointsum B = (A\times
\{0\}) \cup (B \times \{1\})$ છે.
\end{defn}

હવે ક્રમવાચક સંખ્યાઓની જોડીઓ પર એક ક્રમ વ્યાખ્યાયિત
કરીએ:

\begin{defn}
કોઈ પણ ક્રમવાચક સંખ્યાઓ
$\alpha_1, \alpha_2, \beta_1, \beta_2$ માટે કહીએ કે:
\begin{align*}
	\tuple{\alpha_1, \alpha_2} \rlexless \tuple{\beta_1, \beta_2}\text{ iff }&
	\text{either $\alpha_2 \in \beta_2$}\\
	& \text{or both $\alpha_2 = \beta_2$ and $\alpha_1 \in \beta_1$}
\end{align*}
\end{defn}

આ \emph{ઊલટો શબ્દકોશક્રમ} છે, કારણ કે પહેલાં બીજા
ઘટકથી અને પછી પહેલા ઘટકથી ક્રમ નક્કી થાય છે. હવે
યાદ કરો કે આપણે $\alpha \ordplus \beta$ને $\alpha$ની
નકલ પછી $\beta$ની નકલ મૂકવાથી મળતો ક્રમપ્રકાર
ગણવા માગતા હતા. એ માટે કહીએ:

\begin{defn}\ollabel{defordplus}
કોઈ પણ ક્રમવાચક સંખ્યાઓ $\alpha$, $\beta$ માટે
તેમનો સરવાળો $\alpha \ordplus
\beta = \ordtype{\alpha \disjointsum \beta, \rlexless}$ છે.
\end{defn}
\noindent
ધ્યાન આપો કે અહીં સંકેતનો થોડો ઢીલો ઉપયોગ કર્યો છે;
કડક રીતે ``$\rlexless$''ને બદલે
``$\Setabs{\tuple{x,y}\in \alpha\disjointsum\beta}{x \rlexless y}$''
લખવું જોઈએ. જોકે ટૂંકાણ માટે આગળ પણ આ રીતે જ
સંકેત વાપરીશું.

નીચેનું પરિણામ અને
\olref[ordinals][ordtype]{thmOrdinalRepresentation}
મળીને ખાતરી આપે છે કે આપણી વ્યાખ્યા સુગઠિત છે:

\begin{lem}\ollabel{ordsumlessiswo}
કોઈ પણ ક્રમવાચક સંખ્યાઓ $\alpha$ અને $\beta$ માટે
$\tuple{\alpha \disjointsum \beta, \rlexless}$ સુક્રમ છે.
\end{lem}

\begin{proof}
$\alpha \disjointsum \beta$ પર $\rlexless$ સંલગ્ન
છે તે સ્પષ્ટ છે. તે સુઆધારિત છે એમ બતાવવા અરિક્ત
$X \subseteq \alpha
\disjointsum \beta$ નક્કી કરો. $X$માંથી બીજા
સ્થાનનો ઘટક શક્ય એટલો નાનો હોય એવી જોડીઓનો ઉપગણ
$Y$ લો; એટલે કે
$Y = \Setabs{\tuple{\gamma, i} \in X}{(\forall \tuple{\delta, j} \in X)i \leq j}$.
હવે $Y$માંથી પહેલા સ્થાનનો ઘટક સૌથી નાનો હોય એવી
જોડી પસંદ કરો.
\end{proof}
\noindent
આમ ક્રમવાચક સરવાળાની સરસ સ્પષ્ટ વ્યાખ્યા મળી.
હવે એક અપેક્ષિત તથ્ય (\olref[z][infinity-again]{defnomega}
મુજબ $1 = \{0\}$ યાદ કરો):
\begin{prop}
કોઈ પણ ક્રમવાચક સંખ્યા $\alpha$ માટે
$\alpha \ordplus  1 = \ordsucc{\alpha}$.
\end{prop}

\begin{proof}
$\ordsucc{\alpha} = \alpha \cup
\{\alpha\}$થી $\alpha\disjointsum1 = (\alpha \times \{0\})
\cup (\{0\} \times \{1\})$ સુધીની ક્રમ-એકરૂપતા
$f$ લો, જ્યાં $\gamma \in \alpha$ માટે $f(\gamma) =
\tuple{\gamma, 0}$ અને $f(\alpha) = \tuple{0,
1}$.
\sourcecorrection{OLMD-128}{મૂળની આ પંક્તિમાં
પહેલેથી ચિહ્નિત થયેલા બે ગણો વચ્ચે ફરી
``$\disjointsum$'' લખાયું છે; તેથી તે મૂળ
$\alpha\disjointsum1$ સમાન નથી. અહીં વ્યાખ્યા
મુજબ બંને ચિહ્નિત ભાગોનો યોગગણ ``$\cup$''
લખ્યો છે.}
\end{proof}
\noindent
વધુમાં, સરવાળો કેટલીક પુનરાવર્તી શરતો સંતોષે છે તે
બતાવવું સરળ છે:

\begin{lem}\ollabel{ordadditionrecursion}
કોઈ પણ ક્રમવાચક સંખ્યાઓ $\alpha, \beta$ માટે:
\begin{align*}
	\alpha\ordplus 0 &= \alpha\\
	\alpha \ordplus  (\beta\ordplus 1) &= (\alpha \ordplus  \beta) \ordplus  1\\
	\alpha  \ordplus  \beta &= \supstrict_{\delta < \beta}(\alpha \ordplus  \delta) && \text{if $\beta $ is a limit ordinal}
\end{align*}
\end{lem}

\begin{proof}
એક પછી એક કિસ્સો તપાસીએ; પહેલાં:
\begin{align*}
	\alpha \ordplus  0
	& = \ordtype{(\alpha \times \{0\}) \cup (0 \times \{1\}), \rlexless} \\
	&= \ordtype{(\alpha \times \{0\}) \cup \emptyset, \rlexless}\\
	&= \alpha\\
	\alpha \ordplus (\beta \ordplus  1)
	%&= \ordtype{(\alpha\times \{0\}) \cup (\ordtype{\beta \ordplus  1}\times \{1\}), \rlexless} \\
	&= \ordtype{(\alpha\times \{0\}) \cup (\ordsucc{\beta}\times \{1\}), \rlexless} \\
	&= \ordtype{(\alpha\times \{0\}) \cup (\beta \times \{1\}), \rlexless} \ordplus 1\\
	&= (\alpha \ordplus  \beta) \ordplus  1
\end{align*}
\sourcecorrection{OLMD-129}{મૂળમાં $0\times\{1\}$
ખાલી ગણ હોવા છતાં આગળની પંક્તિમાં તે
``$\{0\}$'' બની જાય છે. અહીં તેને
``$\emptyset$'' રાખ્યું છે.}
હવે $\beta \neq \emptyset$ સીમા ક્રમવાચક સંખ્યા
હોય. $\delta < \beta$ હોય તો
$\delta\ordplus 1 < \beta$ પણ છે; તેથી
$\alpha \ordplus  \delta$ એ
$\alpha \ordplus  \beta$નો યોગ્ય આરંભિક ખંડ છે.
આથી $\alpha
\ordplus  \beta$ એ
$X = \Setabs{\alpha
\ordplus  \delta}{\delta < \beta}$ની કડક ઉચ્ચસીમા
છે. વળી, $\alpha \leq \gamma <
\alpha \ordplus  \beta$ હોય, તો સ્પષ્ટપણે કોઈ
$\delta < \beta$ માટે $\gamma = \alpha \ordplus
\delta$. તેથી $\alpha \ordplus  \beta =
\supstrict_{\delta< \beta}(\alpha\ordplus \delta)$.
\end{proof}

અહીં એક નોંધપાત્ર વાત છે. ક્રમવાચક સરવાળો
વ્યાખ્યાયિત કરવા આપણે \emph{તેના બદલે} સીધો
અનંતાતીત પુનરાવર્તનનો પ્રમેય વાપરી શક્યા હોત,
અને \olref{ordadditionrecursion}માં આપેલી પુનરાવર્તી
સમીકરણો જ લઈ શક્યા હોત (જોકે
``$\beta \ordplus 1$''ને બદલે
``$\ordsucc{\beta}$'' વાપરીને).

આમ ક્રમવાચક સંખ્યાઓ પર ક્રિયાઓની વ્યાખ્યા આપવાની
બે જુદી રીતો છે. સ્પષ્ટ સુક્રમિત ગણ રચે અને તેનો
ક્રમપ્રકાર વિચારે, એ રીતે \emph{રચના કરીને} વ્યાખ્યા
આપી શકાય. અથવા માત્ર પુનરાવર્તી સમીકરણો આપીને
\emph{પુનરાવર્તનથી} વ્યાખ્યા આપી શકાય. જોકે યોગ્ય
રીતે કરીએ, તો બંનેના પરિણામ સમાન છે. કારણ કે
\olref[ordinals][ordtype]{thmOrdinalRepresentation}
આ પુનરાવર્તી સમીકરણો \emph{અનન્ય} ક્રમવાચક સંખ્યાઓ
નક્કી કરે છે તેની ખાતરી આપે છે.
\sourcecorrection{OLMD-130}{મૂળમાં પુનરાવર્તી
સમીકરણોની અનન્યતા માટે ક્રમવાચક-રજૂઆત પ્રમેયનો
સંદર્ભ છે. એ પ્રમેય સુક્રમને અનન્ય ક્રમવાચક
ક્રમપ્રકાર આપે છે; પુનરાવર્તી સમીકરણો માટેનું
અસ્તિત્વ અને અનન્યતા અનંતાતીત પુનરાવર્તનના
પ્રમેયથી મળે છે. અહીં મૂળનો સંદર્ભ રાખીને
આ તર્કભેદ નોંધ્યો છે.}

ઘણી રીતે ક્રમવાચક અંકગણિત પ્રાકૃતિક સંખ્યાઓના
સરવાળાની જેમ વર્તે છે. દાખલા તરીકે, નીચેનું સાબિત
કરી શકીએ છીએ:

\begin{lem}\ollabel{ordinaladditionisnice}
$\alpha, \beta, \gamma$ ક્રમવાચક સંખ્યાઓ હોય, તો:
\begin{enumerate}
	\item\ollabel{ordaddition1} જો $\beta < \gamma$ હોય, તો $\alpha
	\ordplus  \beta < \alpha \ordplus  \gamma$
	\item\ollabel{ordaddition2} જો $\alpha \ordplus  \beta =
	\alpha\ordplus \gamma$ હોય, તો $\beta = \gamma$
	\item\ollabel{ordaddition3} $\alpha \ordplus  (\beta \ordplus
	\gamma) = (\alpha \ordplus  \beta) \ordplus  \gamma$;
	એટલે સરવાળો સહયોગી છે
	\item\ollabel{ordaddition4} જો $\alpha \leq \beta$ હોય, તો $\alpha
	\ordplus  \gamma \leq \beta \ordplus \gamma$
\end{enumerate}
\end{lem}

\begin{proof}
અહીં \olref{ordaddition3} સાબિત કરીએ છીએ; બાકીનાં
સ્વાધ્યાય તરીકે છોડીએ છીએ. $\gamma$ પર સરળ
અનંતાતીત અનુમાનથી સાબિતી કરીએ, અને
\olref{ordadditionrecursion} વાપરીએ. $\gamma = 0$
હોય ત્યારે:
\[
(\alpha \ordplus  \beta) \ordplus  0 = \alpha \ordplus  \beta  = \alpha \ordplus  (\beta \ordplus  0)
\]
$\gamma = \delta\ordplus 1$ હોય, ત્યારે અનુમાનની
ધારણા તરીકે $(\alpha
\ordplus  \beta) \ordplus  \delta = \alpha \ordplus  (\beta \ordplus
\delta)$ લો; હવે \olref{ordadditionrecursion} ત્રણ વાર વાપરીએ:
\begin{align*}
	(\alpha \ordplus  \beta) \ordplus  (\delta \ordplus  1) & = ((\alpha \ordplus  \beta) \ordplus  \delta)\ordplus 1\\
	& = (\alpha \ordplus  (\beta \ordplus  \delta)) \ordplus  1\\
	& = \alpha \ordplus  ((\beta \ordplus  \delta)\ordplus 1)\\
	& = \alpha \ordplus  (\beta \ordplus  (\delta\ordplus 1))
\end{align*}
$\gamma$ સીમા ક્રમવાચક સંખ્યા હોય, ત્યારે
$\delta \in \gamma$ હોય તો $(\alpha \ordplus  \beta) \ordplus  \delta =
\alpha \ordplus  (\beta \ordplus  \delta)$ એવી
અનુમાનની ધારણા લો; હવે:
\begin{align*}
	(\alpha \ordplus  \beta) \ordplus  \gamma & = \supstrict_{\delta < \gamma}((\alpha \ordplus  \beta) \ordplus  \delta) \\
	&= \supstrict_{\delta < \gamma}(\alpha \ordplus  (\beta \ordplus  \delta))\\
	&= \alpha \ordplus  \supstrict_{\delta < \gamma}(\beta \ordplus  \delta)\\
	& = \alpha \ordplus  (\beta \ordplus  \gamma)
\end{align*}
\end{proof}

\begin{prob}
\olref[sth][ord-arithmetic][add]{ordinaladditionisnice}નાં
બાકીનાં પરિણામો સાબિત કરો.
\end{prob}

આ રીતે ક્રમવાચક સરવાળો બહુ પરિચિત લાગવો જોઈએ.
પરંતુ એક મહત્ત્વની બાબતમાં તે પ્રાકૃતિક સંખ્યાઓના
સરવાળા જેવો \emph{નથી}.

\begin{prop}\ollabel{ordsumnotcommute}
ક્રમવાચક સરવાળો {અદલબદલીય નથી};
$1 \ordplus  \omega = \omega <
\omega \ordplus  1$.
\end{prop}

\begin{proof}
ધ્યાન આપો કે $1 \ordplus  \omega = \supstrict_{n < \omega} (1 \ordplus n)
= \omega \in \omega \cup \{\omega\} = \ordsucc{\omega} = \omega
\ordplus  1$.
\end{proof}
\noindent
પહેલી વાર આ વાત આશ્ચર્યજનક લાગી શકે, પરંતુ
\emph{એવું લાગવું ન જોઈએ}. એક તરફ,
$1 \ordplus  \omega$ વિચારીએ ત્યારે બધી પ્રાકૃતિક
સંખ્યાઓની \emph{પહેલાં} એક વધારાનું ઘટક મૂકવાથી
મળતો ક્રમપ્રકાર વિચારીએ છીએ.
\olref[sfr][infinite][hilbert]{sec}માં હિલ્બર્ટની હોટેલ
વિશે કરેલી દલીલ મુજબ, સહજ રીતે આ વધારાનું પહેલું
ઘટક આખા ક્રમપ્રકારને બદલવું ન જોઈએ. બીજી તરફ,
$\omega \ordplus  1$ વિચારીએ ત્યારે બધી પ્રાકૃતિક
સંખ્યાઓની \emph{પછી} એક વધારાનું ઘટક મૂકવાથી
મળતો ક્રમપ્રકાર વિચારીએ છીએ. અને તે તો ધરમૂળથી
જુદી બાબત છે!

\end{document}
